\exercisesection{}

\begin{enumerate}
\def\labelenumi{\arabic{enumi})}
\item
  Let G be a locally compact group, H a closed subgroup of G. For \(\xi\in H\) , set \(\chi(\xi)=\Delta_{H}(\xi)\Delta_{G}(\xi)^{-1}\) . Regard H as operating on the right in G by translations. Show that there exists on \(\mathcal{K}^{\chi}(G)\) a nonzero positive linear form I and, up to a constant factor, only one, that is invariant under left translation. (Make use of Prop. 3 with X=G, while taking \(\mu\) to be a left Haar measure of G.)
\item
  Let X and \(X'\) be two locally compact spaces in which a locally compact group H operates on the right, continuously and properly. Let \(\theta\) be a proper continuous mapping of X into \(X'\) , compatible with the identity mapping of H (GT, III, §2, No.~4), and let \(\theta': X/H \to X'/H\) be the mapping deduced from \(\theta\) by passage to quotients. Let \(f'\) be a continuous function on \(X'\) whose support has compact intersection with the saturation of each compact set. Then \(f = f' \circ \theta\) has the same properties in X, and
\end{enumerate}

\[
f ^ {b} = f ^ {\prime b} = \theta^ {\prime}
\]

(the mappings \(f \mapsto f^b\) , \(f' \mapsto f'^b\) being relative to the choice of a same Haar measure on H).

\begin{enumerate}
\def\labelenumi{\arabic{enumi})}
\setcounter{enumi}{2}
\tightlist
\item
  Let B be a locally compact space, H a locally compact group. Set \(\mathbf{X} = \mathbf{B} \times \mathbf{H}\) , the group H operating in X by
\end{enumerate}

\[
(b, \xi) \xi^ {\prime} = (b, \xi \xi^ {\prime}).
\]

Let \(\lambda\) be a measure on \(B = X / H\) , and \(\beta\) a left Haar measure on \(H\) . Then \(\lambda^{\sharp} = \lambda \otimes \beta\) .

\begin{enumerate}
\def\labelenumi{\arabic{enumi})}
\setcounter{enumi}{3}
\item
  Let X be a locally compact space in which a locally compact group H acts on the right continuously and properly. Let \(\beta\) be a left Haar measure on H, \(\mu\) a measure \(\geqslant0\) on X. Let h be a continuous function \(\geqslant0\) on X such that \(h^{b}=1\) . If f is a \(\mu\) -negligible function \(\geqslant0\) on X, the function \((x,\xi)\mapsto f(x)h(x\xi)\) is \((\mu\otimes\beta)\) -negligible on X×H. (There exist decreasing open sets \(\Omega_{1},\Omega_{2},\ldots\) such that \(\mu(\Omega_{n})\) tends to 0 and such that f is zero outside \(\Omega_{1}\cap\Omega_{2}\cap\cdots\) . Show that \(\int^{*}\varphi_{\Omega_{n}}(x)h(x\xi)d\mu(x)d\beta(\xi)\leqslant\mu(\Omega_{n})\) , on observing that the function \((x,\xi)\mapsto\varphi_{\Omega_{n}}(x)h(x\xi)\) is lower semi-continuous.)
\item
  Let \(\mathbf{G}\) be a locally compact group. For every \(s \in \mathbf{G}\) , let \(\psi_s\) be the automorphism of the additive group \(\mathbf{R}\) defined by \(\psi_s(x) = \Delta_{\mathbf{G}}(s)x\) . Let \(\Gamma\) be the topological semi-direct product of \(\mathbf{G}\) and \(\mathbf{R}\) defined by \(s \mapsto \psi_s\) (GT, III, §2, No.~10). Show that \(\Gamma\) is unimodular.
\item
  Let \(G\) be a locally compact group, \(G_1, G_2, G_3\) closed subgroups such that \(G_3 \subset G_1 \cap G_2\) . Assume that \(G, G_1, G_2, G_3\) are unimodular, and that \(G_1 / G_3\) has finite total measure (for every measure invariant under \(G_1\) ). Let \(\lambda, \mu, \nu\) be invariant measures on \(G / G_1, G / G_2, G_2 / G_3\) , and \(\varphi\) the canonical mapping of \(G\) onto \(G / G_1\) . Let \(f \in \mathcal{K}(G / G_1)\) . Show that
\end{enumerate}

\[
a \int_ {\mathrm{G} / \mathrm{G} _ {1}} f (u) d \lambda (u) = \int_ {\mathrm{G} / \mathrm{G} _ {2}} d \mu (\dot {x}) \int_ {\mathrm{G} _ {2} / \mathrm{G} _ {3}} f (\varphi (x \xi)) d \nu (\dot {\xi})
\]

\((\dot{x} = x\mathrm{G}_{2}, \dot{\xi} = \xi \mathrm{G}_{3})\) , where \(a\) is a constant independent of \(f\) .

¶ 7) a) Let E be a locally compact space, Γ a locally compact group operating on the left continuously in E. Assume that for every \(x \in E\) , the mapping \(s \mapsto s \cdot x\) of Γ into E is proper, and that there exists a nonzero bounded positive measure μ on E invariant under Γ. Then Γ is compact. (Let \((s_n)\) be a sequence of points of Γ. Let K be a compact subset of E that is not μ-negligible. Show, using Exer. 15 of Ch. IV, §4, that there exists an \(x_0 \in E\) such that \(s_n x_0 \in K\) for infinitely many values of n.~From this, deduce that \((s_n)\) has a cluster point in Γ. Then apply Exer. 6 of GT, II, §4.)

\begin{enumerate}
\def\labelenumi{\alph{enumi})}
\setcounter{enumi}{1}
\tightlist
\item
  Let G be a locally compact group, H and K two closed subgroups of G. Assume that G and H are unimodular and that G/H has finite measure for the measure invariant under G. Show that for H and K to satisfy the equivalent conditions of Exer. 11 c) of GT, III, §4, it is necessary and sufficient that K be compact.
\end{enumerate}

\begin{enumerate}
\def\labelenumi{\arabic{enumi})}
\setcounter{enumi}{7}
\tightlist
\item
  Let \(\mathfrak{G}\) be a locally compact group, G and H two closed subgroups of \(\mathfrak{G}\) . Assume that every \(s \in \mathfrak{G}\) can be written in one and only one way in the form \(s = \xi x = y\eta\) \((\xi, \eta \in \mathbf{G}, x, y \in \mathbf{H})\) , and that \(\xi, x, y, \eta\) depend continuously on \(s\) . Every \(\xi \in \mathbf{G}\) defines a homeomorphism \(\widehat{\xi}\) of H onto H by \(\xi x \in \widehat{\xi}(x)\mathbf{G} (x \in \mathbf{H})\) . Every \(x \in \mathbf{H}\) defines a homeomorphism \(\widehat{x}\) of G onto G by \(\xi x \in \mathrm{H}\widehat{x}(\xi) (\xi \in \mathbf{G})\) . Let \(\mu, \alpha, \beta\) be left Haar measures on \(\mathfrak{G}, \mathbf{G}, \mathbf{H}\) . Show that
\end{enumerate}

\[
\begin{array}{c} d \widehat {x} ^ {- 1} (\alpha) (\xi) = \frac {\Delta_ {\mathfrak {G}} (\widehat {\xi} (x)) \Delta_ {\mathrm{G}} (\widehat {x} (\xi))}{\Delta_ {\mathrm{H}} (\widehat {\xi} (x)) \Delta_ {\mathrm{G}} (\xi)} d \alpha (\xi), \\ d \widehat {\xi} ^ {- 1} (\beta) (x) = \frac {\Delta_ {\mathrm{G}} (\widehat {x} (\xi))}{\Delta_ {\mathfrak {G}} (\widehat {x} (\xi))} d \beta (x). \end{array}
\]

(Let \(f \in \mathcal{K}(G)\) . Express the fact that \(\int f(u) d\mu(u)\) , calculated with the help of (23), either for \(X = G\) , \(Y = H\) , or for \(X = H\) , \(Y = G\) , is invariant when \(f\) is subjected to left translation by an element of \(G\) or an element of \(H\) .)

\begin{enumerate}
\def\labelenumi{\arabic{enumi})}
\setcounter{enumi}{8}
\tightlist
\item
  Let X be a locally compact space, H a compact group operating continuously (hence properly) on the right in X. Denote by \(\beta\) the normalized Haar measure on H, by \(\pi: X \to X/H\) the canonical mapping. Let \(\lambda\) be a positive measure on X/H, \(\lambda^{\sharp}\) the corresponding measure on X.
\end{enumerate}

\begin{enumerate}
\def\labelenumi{\alph{enumi})}
\item
  Show that if N is a \(\lambda\) -negligible subset of X/H, then \(\overline{\pi}^{1}(N)\) is \(\lambda^{\sharp}\) -negligible. (Calculate the measure of a set \(\overline{\pi}^{1}(U)\) , where U is open in X/H.)
\item
  Let \(p\) be a finite number \(\geqslant 1\) , and let \(f\) be a function in \(\mathcal{L}_{\mathrm{F}}^{p}(\mathrm{X},\lambda^{\sharp})\) (F a Banach space or \(\mathrm{F} = \overline{\mathbf{R}}\) ). Show that the set of \(x \in \mathrm{X}\) such that \(\xi \mapsto f(x\xi)\) is not \(\beta\) -integrable on \(\mathrm{H}\) is of the form \(\overline{\pi}^{1}(\mathrm{N})\) , where \(\mathrm{N}\) is \(\lambda\) -negligible. Moreover, if \(f^b\) is the function on \(\mathrm{X / H}\) , defined almost everywhere (for \(\lambda\) ), such that \(f^b (\pi (x)) = \int_{\mathrm{H}} f(x\xi) d\beta (\xi)\) , then \(f^b\) belongs to \(\mathcal{L}_{\mathrm{F}}^{p}(\mathrm{X / H},\lambda)\) , and \(\mathrm{N}_p(f^b)\leqslant \mathrm{N}_p(f)\) .
\item
  Conversely, if \(g \in \mathcal{L}_{\mathrm{F}}^{p}(\mathrm{X / H}, \lambda)\) , then the function \(g \circ \pi\) belongs to \(\mathcal{L}_{\mathrm{F}}^{p}(\mathrm{X}, \lambda^{\sharp})\) . If \(p, q\) are conjugate exponents, \(\mathrm{F}'\) is the dual of \(\mathrm{F}\) , \(f \in \mathcal{L}_{\mathrm{F}}^{p}(\mathrm{X}, \lambda^{\sharp})\) and \(g \in \mathcal{L}_{\mathrm{F}'}^{q}(\mathrm{X / H}, \lambda)\) , then
\end{enumerate}

\[
\int_ {\mathrm{X}} \left\langle f (x), g (\pi (x)) \right\rangle d \lambda^ {\sharp} (x) = \int_ {\mathrm{X/H}} \left\langle f ^ {\flat} (z), g (z) \right\rangle d \lambda (z).
\]

\begin{enumerate}
\def\labelenumi{\arabic{enumi})}
\setcounter{enumi}{9}
\tightlist
\item
  With the notations of §1, No.~6, Prop. 6 let, for every \(\alpha \in A\) , \(\mu_{\alpha}\) be a left Haar measure on \(G_{\alpha}\) , so that the inverse limit of the \(\mu_{\alpha}\) is a left Haar measure \(\mu\) on \(G\) . For every \(\alpha \in A\) , denote by \(\lambda_{\alpha}\) the normalized Haar measure on \(K_{\alpha}\) .
\end{enumerate}

\begin{enumerate}
\def\labelenumi{\alph{enumi})}
\tightlist
\item
  Let \(p\) be a finite number \(\geqslant 1\) and let \(f\) be a function in \(\mathcal{L}_{\mathrm{F}}^{p}(\mathrm{G},\mu)\) (F a Banach space or \(\mathrm{F} = \overline{\mathbf{R}}\) ). For every \(\alpha \in \mathbf{A}\) , the function
\end{enumerate}

\[
f _ {\alpha} (s) = \int_ {K _ {\alpha}} f (s \xi) d \lambda_ {\alpha} (\xi),
\]

defined almost everywhere for \(\mu\) , belongs to \(\mathcal{L}_{\mathrm{F}}^{p}(\mathrm{G},\mu)\) (Exer. 9). Show that, with respect to the directed set A, \(f_{\alpha}\) tends in mean of order \(p\) to \(f\) (make use of Lemma 2 of §1, No.~6).

\begin{enumerate}
\def\labelenumi{\alph{enumi})}
\setcounter{enumi}{1}
\tightlist
\item
  Assume that A = N and p = 1. Show that \(f_{n}\) then tends to f almost everywhere (for \(\mu\) ) in G (use a method analogous to that of Ch. V, §8, Exer. 16).
\end{enumerate}

¶ 11) Let G be a locally compact group, H a closed subgroup of G, λ a left Haar measure on G; assume that there exists on G/H a measure μ invariant under G such that λ = μ \(^{\sharp}\) and μ(G/H) \textless{} +∞. Let ν be a nonzero positive measure on G/H such that \(\nu (\mathrm{G / H}) < +\infty\) . Finally, let \(h\) be a function on \(\mathbf{G}\) having the properties of No.~4, Prop. 8.

\begin{enumerate}
\def\labelenumi{\alph{enumi})}
\tightlist
\item
  Let A be a Borel subset of G/H. Show that
\end{enumerate}

\[
\int_ {\mathrm{G}} h (s) \nu (s ^ {- 1} \mathrm{A}) d \lambda (s) = \nu (\mathrm{G/H}) \mu (\mathrm{A})
\]

(make use of Prop. 9 a) of No.~4).

\begin{enumerate}
\def\labelenumi{\alph{enumi})}
\setcounter{enumi}{1}
\tightlist
\item
  Let \(A_i\) ( \(1 \leqslant i \leqslant n\) ) be Borel subsets of \(G / H\) and, for every \(i\) , let \(b_i = \sup_{s \in G} \nu(s^{-1} A_i)\) . Show that there exist two distinct indices \(i, j\) such that
\end{enumerate}

\[
\int_ {\mathrm{G}} h (s) \nu (s ^ {- 1} \mathrm{A} _ {i}) \nu (s ^ {- 1} \mathrm{A} _ {j}) d \lambda (s) \geqslant \frac {\left(\nu (\mathrm{G} / \mathrm{H})\right) ^ {2}}{\mu (\mathrm{G} / \mathrm{H})} \left(\left(\sum_ {k = 1} ^ {n} \mu (\mathrm{A} _ {k})\right) ^ {2} - \frac {\mu (\mathrm{G} / \mathrm{H})}{\nu (\mathrm{G} / \mathrm{H})} \sum_ {k = 1} ^ {n} b _ {k} \mu (\mathrm{A} _ {k})\right)
\]

(argue as in §1, Exer. 28 a)).

¶ 12) a) Let X be a topological space, \(f: X \to N\) an upper semi-continuous function. Let \(X_0\) be the set of points of X where \(f\) is locally constant, and \(Y_0 = X - X_0\) . Define recursively two sequences \((X_i)_{i \geqslant 0}\) , \((Y_i)_{i \geqslant 0}\) of subsets of X, in the following way: \(X_i\) is the set of points of \(Y_{i-1}\) where \(f|_{Y_{i-1}}\) is locally constant, and \(Y_i = Y_{i-1} - X_i\) . Show that \(X_i\) is a dense open subset of \(Y_{i-1}\) , and that \(\bigcap_i Y_i = \varnothing\) .

(Show that \(f(x) > i\) at every point of \(Y_i\) .)

\begin{enumerate}
\def\labelenumi{\alph{enumi})}
\setcounter{enumi}{1}
\tightlist
\item
  Let X be a set, H a group operating on the right in X, \(\pi\) the canonical mapping of X onto X/H, and A (resp. B) the set of A ⊂ X such that \(\pi|A:A\to X/H\) is injective (resp. surjective). Let \((V_{1},V_{2},\ldots)\) be a countable covering of X by elements of A. Let
\end{enumerate}

\[
\mathrm{V} _ {i} ^ {\prime} = \mathrm{V} _ {i} \cap \complement \left(\mathrm{V} _ {1} \mathrm{H} \cup \mathrm{V} _ {2} \mathrm{H} \cup \dots \cup \mathrm{V} _ {i - 1} \mathrm{H}\right).
\]

Show that \(\mathrm{F} = \bigcup_{i\geqslant 0}\mathrm{V}_i^\prime \in \mathfrak{A}\cap \mathfrak{B}\) .

\begin{enumerate}
\def\labelenumi{\alph{enumi})}
\setcounter{enumi}{2}
\item
  Let X be a locally compact space, H a countable discrete group operating on the right in X continuously and properly, \(\pi\) the canonical mapping of X onto X/H, and \(\mu\) a measure \(\geqslant 0\) on X. Show that if the sets \(V_{i}\) of b) are quadrable (Ch. IV, §5, Exer. 17 d)), then the set F of b) is a quadrable fundamental domain.
\item
  Let us maintain the hypotheses of \(c\) ), and use the notations \(\mathbf{H}_x, n(x)\) of No.~10. An \(x \in X\) is said to be general if there exists a neighborhood \(V\) of \(x\) in \(X\) such that \(\mathrm{H}_y = \mathrm{H}_x\) for all \(y \in V\) . Show that the general points of \(X\) are those where the function \(n\) is locally constant. Show that a general point admits an open neighborhood that belongs to \(\mathfrak{A}\) and is quadrable (make use of Ch. IV, §5, Exer. 17 d)).
\item
  Let us maintain the hypotheses of \(c\) ), and assume moreover that \(X\) is countable at infinity. Show that if \(X - X_0\) is \(\mu\) -negligible, then there exists a fundamental domain \(F\) that is Borel and quadrable. (Apply the construction of \(a\) ) to the function \(n\) . Then apply the results of \(c\) ) and \(d\) ) in \(X_0\) . Argue similarly in each \(X_i\) .)
\item
  Let \(U \in B\) be an open set. Show that one can impose on the set F of e) the following supplementary properties: 1) \(F \subset U\) ; 2) for every compact subset K of X, the set of \(s \in H\) such that Fs intersects K is finite.
\end{enumerate}

\begin{enumerate}
\def\labelenumi{\arabic{enumi})}
\setcounter{enumi}{12}
\tightlist
\item
  Let G be a compact group, \(\mu\) a Haar measure on G, u an endomorphism of G such that \(u(\mathrm{G})\) is an open subgroup of G and the kernel \(\overline{u}(e)\) (denoted \(G_{u}\) ) a finite subgroup of G.
\end{enumerate}

\begin{enumerate}
\def\labelenumi{\alph{enumi})}
\item
  Show that there exist a real number \(h(u) > 0\) and an open neighborhood \(U\) of \(e\) in \(G\) such that, for every open set \(V \subset U\) , \(u(V)\) is open in \(G\) and \(\mu(u(V)) = h(u)\mu(V)\) (cf.~§1, No.~7, Cor. of Prop. 9).
\item
  Show that \(h(u) = \operatorname{Card}(\mathrm{G} / u(\mathrm{G})) / \operatorname{Card}(\mathrm{G}_u)\) (calculate \(\mu(u(\mathrm{G}))\) in two ways, using a) and Prop. 10 of No.~7).
\end{enumerate}

